\documentclass[../main.tex]{subfiles}

\begin{document}

\chapter{泛函分析}\label{chapter:appD-泛函分析}

\section{Banach 空间}\label{section:appD-Banach空间}

设 $X$ 表示实线性空间。

\begin{definition}\label{definition:appD-范数}
映射 $\|\cdot\|:X\to[0,\infty)$ 称为范数，若
\begin{enumerate}[(i)]
\item $\|u+v\|\leqslant\|u\|+\|v\|$ 对所有 $u,v\in X$，
\item $\|\lambda u\|=|\lambda|\|u\|$ 对所有 $u\in X$，$\lambda\in\mathbb{R}$，
\item $\|u\|=0$ 当且仅当 $u=0$。
\end{enumerate}
不等式 (i) 是三角不等式。
\end{definition}

此后假设 $X$ 是赋范线性空间。

\begin{definition}\label{definition:appD-收敛}
称序列 $\{u_k\}_{k=1}^\infty\subset X$ 收敛到 $u\in X$，写 $u_k\to u$，若
\begin{align*}
\lim_{k\to\infty}\|u_k-u\|=0.
\end{align*}
\end{definition}

\begin{definition}\label{definition:appD-完备}
\begin{enumerate}[(i)]
\item 序列 $\{u_k\}_{k=1}^\infty\subset X$ 称为 Cauchy 序列，只要对每个 $\varepsilon>0$ 存在 $N>0$ 使得
\begin{align*}
\|u_k-u_l\|<\varepsilon
\end{align*}
对所有 $k,l\geqslant N$。

\item $X$ 是完备的，若 $X$ 中每个 Cauchy 序列收敛；即每当 $\{u_k\}_{k=1}^\infty$ 是 Cauchy 序列，存在 $u\in X$ 使得 $\{u_k\}_{k=1}^\infty$ 收敛到 $u$。

\item Banach 空间 $X$ 是完备的赋范线性空间。
\end{enumerate}
\end{definition}

\begin{definition}\label{definition:appD-可分}
称 $X$ 可分，若 $X$ 包含可数稠密子集。
\end{definition}

\begin{remark}\label{remark:appD-例子}
\textbf{例子。}
\begin{enumerate}[(i)]
\item $L^p$ 空间。假设 $U$ 是 $\mathbb{R}^n$ 的开子集，$1\leqslant p\leqslant\infty$。若 $f:U\to\mathbb{R}$ 可测，定义
\begin{align*}
\|f\|_{L^p(U)}\coloneqq
\begin{cases}
\left(\int_U|f|^p\,\mathrm{d}x\right)^{1/p} & \text{若 } 1\leqslant p<\infty,\\
\operatorname*{ess\,sup}_U|f| & \text{若 } p=\infty.
\end{cases}
\end{align*}
定义 $L^p(U)$ 为所有可测函数 $f:U\to\mathbb{R}$（$\|f\|_{L^p(U)}<\infty$）的线性空间。则 $L^p(U)$ 是 Banach 空间，只要我们将 a.e. 相等的两个函数等同。

\item H\"older 空间。见 \S5.1。

\item Sobolev 空间。见 \S5.2。
\end{enumerate}
\end{remark}

\section{Hilbert 空间}\label{section:appD-Hilbert空间}

设 $H$ 是实线性空间。

\begin{definition}\label{definition:appD-内积}
映射 $(,):H\times H\to\mathbb{R}$ 称为内积，若
\begin{enumerate}[(i)]
\item $(u,v)=(v,u)$ 对所有 $u,v\in H$，
\item 映射 $u\mapsto(u,v)$ 对每个 $v\in H$ 是线性的，
\item $(u,u)\geqslant 0$ 对所有 $u\in H$，
\item $(u,u)=0$ 当且仅当 $u=0$。
\end{enumerate}
\end{definition}

\begin{definition}\label{definition:appD-Hilbert范数}
若 $(,)$ 是内积，关联范数为
\begin{align}
\|u\|\coloneqq(u,u)^{1/2} \quad (u\in H).
\label{align:appD-Hilbert范数}
\end{align}
Cauchy--Schwarz 不等式陈述
\begin{align}
|(u,v)|\leqslant\|u\|\|v\| \quad (u,v\in H).
\label{align:appD-HilbertCS}
\end{align}
此不等式如 \S B.2 中证明。利用\eqref{align:appD-HilbertCS}，容易验证\eqref{align:appD-Hilbert范数} 定义了 $H$ 上的范数。
\end{definition}

\begin{definition}\label{definition:appD-Hilbert空间}
Hilbert 空间 $H$ 是赋予内积生成范数的 Banach 空间。
\end{definition}

\begin{remark}\label{remark:appD-Hilbert例子}
\textbf{例子。}
a. 空间 $L^2(U)$ 是 Hilbert 空间，内积
\begin{align*}
(f,g)=\int_U fg\,\mathrm{d}x.
\end{align*}
b. Sobolev 空间 $H^1(U)$ 是 Hilbert 空间，内积
\begin{align*}
(f,g)=\int_U fg+Df\cdot Dg\,\mathrm{d}x.
\end{align*}
\end{remark}

\begin{definition}\label{definition:appD-正交}
\begin{enumerate}[(i)]
\item 两个元素 $u,v\in H$ 正交，若 $(u,v)=0$。

\item 可数基 $\{w_k\}_{k=1}^\infty\subset H$ 称为正交归一的，若
\begin{align*}
\begin{cases}
(w_k,w_l)=0 & (k,l=1,\cdots;\;k\neq l),\\
\|w_k\|=1 & (k=1,\cdots).
\end{cases}
\end{align*}
若 $u\in H$ 且 $\{w_k\}_{k=1}^\infty\subset H$ 是正交归一基，可写
\begin{align*}
u=\sum_{k=1}^\infty(u,w_k)w_k,
\end{align*}
级数在 $H$ 中收敛。此外
\begin{align*}
\|u\|^2=\sum_{k=1}^\infty(u,w_k)^2.
\end{align*}

\item 若 $S$ 是 $H$ 的子空间，$S^\perp=\{u\in H\mid(u,v)=0\text{ 对所有 }v\in S\}$ 是正交于 $S$ 的子空间。
\end{enumerate}
\end{definition}

\section{有界线性算子}\label{section:appD-有界线性算子}

\subsection{Banach 空间上的线性算子}\label{subsection:appD-Banach算子}

设 $X$ 和 $Y$ 是实 Banach 空间。

\begin{definition}\label{definition:appD-线性算子}
\begin{enumerate}[(i)]
\item 映射 $A:X\to Y$ 是线性算子，只要
\begin{align*}
A[\lambda u+\mu v]=\lambda Au+\mu Av
\end{align*}
对所有 $u,v\in X$，$\lambda,\mu\in\mathbb{R}$。

\item $A$ 的值域为 $R(A)\coloneqq\{v\in Y\mid v=Au\text{ 某 }u\in X\}$，$A$ 的零空间为 $N(A)\coloneqq\{u\in X\mid Au=0\}$。
\end{enumerate}
\end{definition}

\begin{definition}\label{definition:appD-有界算子}
线性算子 $A:X\to Y$ 有界，若
\begin{align*}
\|A\|\coloneqq\sup\{\|Au\|_Y\mid\|u\|_X\leqslant 1\}<\infty.
\end{align*}
容易验证有界线性算子 $A:X\to Y$ 是连续的。
\end{definition}

\begin{definition}\label{definition:appD-闭算子}
线性算子 $A:X\to Y$ 称为闭的，只要 $u_k\to u$ 在 $X$ 中且 $Au_k\to v$ 在 $Y$ 中，则 $Au=v$。
\end{definition}

\begin{theorem}[闭图像定理。]\label{theorem:appD-闭图像}设 $A:X\to Y$ 是闭线性算子。则 $A$ 有界。
\end{theorem}

\begin{definition}\label{definition:appD-谱}
设 $A:X\to X$ 是有界线性算子。
\begin{enumerate}[(i)]
\item $A$ 的预解集为
\begin{align*}
\rho(A)=\{\eta\in\mathbb{R}\mid(A-\eta I)\text{ 是一一对应且到上的}\}.
\end{align*}
\item $A$ 的谱为
\begin{align*}
\sigma(A)=\mathbb{R}\setminus\rho(A).
\end{align*}
若 $\eta\in\rho(A)$，闭图像定理蕴含逆 $(A-\eta I)^{-1}:X\to X$ 是有界线性算子。
\end{enumerate}
\end{definition}

\begin{definition}\label{definition:appD-特征值}
\begin{enumerate}[(i)]
\item 称 $\eta\in\sigma(A)$ 是 $A$ 的特征值，只要
\begin{align*}
N(A-\eta I)\neq\{0\}.
\end{align*}
写 $\sigma_p(A)$ 表示 $A$ 的特征值集合；$\sigma_p(A)$ 是点谱。

\item 若 $\eta$ 是特征值且 $w\neq 0$ 满足
\begin{align*}
Aw=\eta w,
\end{align*}
称 $w$ 是关联的特征向量。
\end{enumerate}
\end{definition}

\begin{definition}\label{definition:appD-对偶}
\begin{enumerate}[(i)]
\item 有界线性算子 $u^*:X\to\mathbb{R}$ 称为 $X$ 上的有界线性泛函。

\item 写 $X^*$ 表示 $X$ 上所有有界线性泛函的集合；$X^*$ 是 $X$ 的对偶空间。

\item 若 $u\in X$，$u^*\in X^*$，写
\begin{align*}
\langle u^*,u\rangle
\end{align*}
表示实数 $u^*(u)$。符号 $\langle,\rangle$ 表示 $X^*$ 和 $X$ 的配对。

\item 定义
\begin{align*}
\|u^*\|\coloneqq\sup\{\langle u^*,u\rangle\mid\|u\|\leqslant 1\}.
\end{align*}

\item Banach 空间是自反的，若 $(X^*)^*=X$。更精确地，这意味着对每个 $u^{**}\in(X^*)^*$，存在 $u\in X$ 使得
\begin{align*}
\langle u^{**},u^*\rangle=\langle u^*,u\rangle
\end{align*}
对所有 $u^*\in X^*$。
\end{enumerate}
\end{definition}

\subsection{Hilbert 空间上的线性算子}\label{subsection:appD-Hilbert算子}

现设 $H$ 是实 Hilbert 空间，内积为 $(,)$。

\begin{theorem}[Riesz 表示定理。]\label{theorem:appD-Riesz}$H^*$ 可典范等同于 $H$；更精确地，对每个 $u^*\in H^*$ 存在唯一元素 $u\in H$ 使得
\begin{align*}
\langle u^*,v\rangle=(u,v)
\end{align*}
对所有 $v\in H$。
\end{theorem}

映射 $u^*\mapsto u$ 是 $H^*$ 到 $H$ 的线性同构。

\begin{definition}\label{definition:appD-伴随}
\begin{enumerate}[(i)]
\item 若 $A:H\to H$ 是有界线性算子，其伴随 $A^*:H\to H$ 满足
\begin{align*}
(Au,v)=(u,A^*v)
\end{align*}
对所有 $u,v\in H$。

\item $A$ 对称，若 $A^*=A$。
\end{enumerate}
\end{definition}

\section{弱收敛}\label{section:appD-弱收敛}

设 $X$ 表示实 Banach 空间。

\begin{definition}\label{definition:appD-弱收敛}
称序列 $\{u_k\}_{k=1}^\infty\subset X$ 弱收敛到 $u\in X$，写
\begin{align*}
u_k\rightharpoonup u,
\end{align*}
若
\begin{align*}
\langle u^*,u_k\rangle\to\langle u^*,u\rangle
\end{align*}
对每个有界线性泛函 $u^*\in X^*$。
\end{definition}

容易验证若 $u_k\to u$，则 $u_k\rightharpoonup u$。任何弱收敛序列是有界的。此外，若 $u_k\rightharpoonup u$，则
\begin{align*}
\|u\|\leqslant\liminf_{k\to\infty}\|u_k\|.
\end{align*}

\begin{theorem}[弱紧性。]\label{theorem:appD-弱紧}设 $X$ 是自反 Banach 空间，假设序列 $\{u_k\}_{k=1}^\infty\subset X$ 有界。则存在子序列 $\{u_{k_j}\}_{j=1}^\infty\subset\{u_k\}_{k=1}^\infty$ 和 $u\in X$ 使得
\begin{align*}
u_{k_j}\rightharpoonup u.
\end{align*}
\end{theorem}

换言之，自反 Banach 空间中的有界序列是弱预紧的。特别地，Hilbert 空间中的有界序列包含弱收敛子序列。

Mazur 定理断言 $X$ 的凸闭子集是弱闭的。

\begin{remark}\label{remark:appD-弱收敛例子}
重要例子。我们最常在以下情形使用弱收敛概念。取 $U\subset\mathbb{R}^n$ 开，假设 $1\leqslant p<\infty$。则 $X=L^p(U)$ 的对偶空间是 $X^*=L^q(U)$，其中 $\dfrac{1}{p}+\dfrac{1}{q}=1$，$1<q\leqslant\infty$。更精确地，$L^p(U)$ 上每个有界线性泛函可表示为 $f\mapsto\int_U gf\,\mathrm{d}x$ 对某 $g\in L^q(U)$。因此
\begin{align*}
f_k\rightharpoonup f \quad \text{弱在 } L^p(U) \text{ 中}
\end{align*}
意味着
\begin{align}
\int_U gf_k\,\mathrm{d}x\to\int_U gf\,\mathrm{d}x
\label{align:appD-弱收敛Lp}
\end{align}
当 $k\to\infty$，对所有 $g\in L^q(U)$。

$L^q(U)$ 作为 $L^p(U)$ 对偶的等同表明 $L^p(U)$ 是自反的，若 $1<p<\infty$。特别地，\cref{theorem:appD-弱紧} 保证从 $L^p(U)$（$1<p<\infty$）中的有界序列可提取弱收敛子序列，即满足\eqref{align:appD-弱收敛Lp} 的序列。这是重要的紧性断言，但需非常仔细地注意：收敛\eqref{align:appD-弱收敛Lp} 不蕴含 $f_k\to f$ 逐点或几乎处处。例如，函数 $\{f_k\}_{k=1}^\infty$ 完全可能当 $k\to\infty$ 时振荡越来越快。（见\cref{exercise:ch08-7-1} 和\cref{exercise:ch09-7-2}。）
\end{remark}

\section{紧算子、Fredholm 理论}\label{section:appD-紧算子}

设 $X$ 和 $Y$ 是实 Banach 空间。

\begin{definition}\label{definition:appD-紧算子}
有界线性算子
\begin{align*}
K:X\to Y
\end{align*}
称为紧的，只要对每个有界序列 $\{u_k\}_{k=1}^\infty\subset X$，序列 $\{Ku_k\}_{k=1}^\infty$ 在 $Y$ 中预紧；即存在子序列 $\{u_{k_j}\}_{j=1}^\infty$ 使得 $\{Ku_{k_j}\}_{j=1}^\infty$ 在 $Y$ 中收敛。
\end{definition}

现设 $H$ 表示实 Hilbert 空间，内积为 $(,)$。容易看到若线性算子 $K:H\to H$ 紧且 $u_k\rightharpoonup u$，则 $Ku_k\to Ku$。

\begin{theorem}[伴随的紧性。]\label{theorem:appD-伴随紧}若 $K:H\to H$ 紧，则 $K^*:H\to H$ 也紧。
\end{theorem}

\begin{proof}
设 $\{u_k\}_{k=1}^\infty$ 是 $H$ 中有界序列，提取弱收敛子序列 $u_{k_j}\rightharpoonup u$ 在 $H$ 中。将证明 $K^*u_{k_j}\to K^*u$。事实上，
\begin{align*}
\|K^*u_{k_j}-K^*u\|^2=(K^*u_{k_j}-K^*u,K^*[u_{k_j}-u])=(KK^*u_{k_j}-KK^*u,u_{k_j}-u).
\end{align*}
由于 $K^*$ 线性，$K^*u_{k_j}\rightharpoonup K^*u$，故 $KK^*u_{k_j}\to KK^*u$。因此 $K^*u_{k_j}\to K^*u$。
\end{proof}

\begin{theorem}[Fredholm 二择一。]\label{theorem:appD-Fredholm}设 $K:H\to H$ 是紧线性算子。则
\begin{enumerate}[(i)]
\item $N(I-K)$ 有限维，
\item $R(I-K)$ 闭，
\item $R(I-K)=N(I-K^*)^\perp$，
\item $N(I-K)=\{0\}$ 当且仅当 $R(I-K)=H$，
\item $\dim N(I-K)=\dim N(I-K^*)$。
\end{enumerate}
\end{theorem}

\begin{proof}
1. 若 $\dim N(I-K)=+\infty$，可找到无穷正交规范集 $\{u_k\}_{k=1}^\infty\subset N(I-K)$。则 $Ku_k=u_k$（$k=1,\cdots$）。现在 $\|u_k-u_l\|^2=\|u_k\|^2-2(u_k,u_l)+\|u_l\|^2=2$（$k\neq l$），故 $\|Ku_k-Ku_l\|=\sqrt{2}$（$k\neq l$）。然而这与 $K$ 的紧性矛盾，因 $\{Ku_k\}_{k=1}^\infty$ 不含收敛子序列。断言 (i) 得证。

2. 接下来断言存在常数 $\gamma>0$ 使得
\begin{align}
\|u-Ku\|\geqslant\gamma\|u\|
\label{align:appD-Fredholm估计}
\end{align}
对所有 $u\in N(I-K)^\perp$。事实上，若不然，对 $k=1,\cdots$ 存在元素 $u_k\in N(I-K)^\perp$，$\|u_k\|=1$ 且 $\|u_k-Ku_k\|<1/k$。因此
\begin{align}
u_k-Ku_k\to 0.
\label{align:appD-uK}
\end{align}
但因 $\{u_k\}_{k=1}^\infty$ 有界，存在弱收敛子序列 $u_{k_j}\rightharpoonup u$。由紧性 $Ku_{k_j}\to Ku$，\eqref{align:appD-uK} 蕴含 $u_{k_j}\to Ku=u$。因此 $u\in N(I-K)$，故
\begin{align*}
(u_{k_j},u)=0 \quad (j=1,\cdots).
\end{align*}
令 $k_j\to\infty$ 导出矛盾。

3. 接下来设 $\{v_k\}_{k=1}^\infty\subset R(I-K)$，$v_k\to v$。可找到 $u_k\in N(I-K)^\perp$ 解 $u_k-Ku_k=v_k$。利用\eqref{align:appD-Fredholm估计}，推出
\begin{align*}
\|v_k-v_l\|\geqslant\gamma\|u_k-u_l\|.
\end{align*}
因此 $u_k\to u$ 且 $u-Ku=v$。这证明了 (ii)。

4. 断言 (iii) 现在是 (ii) 和一般事实 $R(A)=N(A^*)^\perp$（对每个有界线性算子 $A:H\to H$）的推论。

5. 为验证 (iv)，先假设 $N(I-K)=\{0\}$，但 $H_1=(I-K)(H)\subsetneq H$。由 (ii)，$H_1$ 是 $H$ 的闭子空间。此外 $H_2\equiv(I-K)(H_1)\subsetneq H_1$，因 $I-K$ 是一一的。类似地写 $H_k\equiv(I-K)^k(H)$（$k=1,\cdots$），看到 $H_k$ 是 $H$ 的闭子空间，$H_{k+1}\subsetneq H_k$（$k=1,\cdots$）。

选取 $u_k\in H_k$，$\|u_k\|=1$，$u_k\in H_{k+1}^\perp$。则 $Ku_k-Ku_l=-(u_k-Ku_k)+(u_l-Ku_l)+(u_k-u_l)$。若 $k>l$，$H_{k+1}\subsetneq H_k\subseteq H_{l+1}\subsetneq H_l$。因此 $u_k-Ku_k$，$u_l-Ku_l$，$u_k\in H_{l+1}$。因 $u_l\in H_{l+1}^\perp$，$\|u_l\|=1$，推出
\begin{align*}
\|Ku_k-Ku_l\|\geqslant 1 \quad (k,l=1,\cdots).
\end{align*}
但这不可能，因 $K$ 紧。

6. 现反过来假设 $R(I-K)=H$。则由 (iii)，$N(I-K^*)=\{0\}$。因 $K^*$ 紧，可利用步骤 5 得出 $R(I-K^*)=H$。但此时 $N(I-K)=R(I-K^*)^\perp=\{0\}$。此结论和步骤 5 完成断言 (iv) 的证明。

7. 接下来断言
\begin{align*}
\dim N(I-K)\geqslant\dim R(I-K)^\perp.
\end{align*}
为此反设 $\dim N(I-K)<\dim R(I-K)^\perp$。则存在有界线性映射 $A:N(I-K)\to R(I-K)^\perp$ 是一一的但不到上的。将 $A$ 延拓为线性映射 $A:H\to R(I-K)^\perp$，令 $Au=0$（$u\in N(I-K)^\perp$）。现在 $A$ 有有限维值域，故 $A$ 从而 $K+A$ 紧。此外 $N(I-(K+A))=\{0\}$。事实上，若 $Ku+Au=u$，则 $u-Ku=Au\in R(I-K)^\perp$。因此 $u-Ku=Au=0$。故 $u\in N(I-K)$，实际上 $u=0$，因 $A$ 在 $N(I-K)$ 上一一。现对 $\widetilde{K}=K+A$ 应用断言 (iv)。得出 $R(I-(K+A))=H$。但这不可能：若 $v\in R(I-K)^\perp$ 但 $v\notin R(A)$，方程
\begin{align*}
u-(Ku+Au)=v
\end{align*}
无解。

8. 由于 $R(I-K^*)^\perp=N(I-K)$，由步骤 7 推出
\begin{align*}
\dim N(I-K^*)\geqslant\dim R(I-K^*)^\perp=\dim N(I-K).
\end{align*}
反向不等式通过交换 $K$ 和 $K^*$ 的角色得到。这建立了 (v)。
\end{proof}

\begin{remark}\label{remark:appD-Fredholm二择一}
\cref{theorem:appD-Fredholm} 特别断言要么
\begin{enumerate}[($\alpha$)]
\item 对每个 $f\in H$，方程 $u-Ku=f$ 有唯一解，
\end{enumerate}
要么
\begin{enumerate}[($\beta$)]
\item 齐次方程 $u-Ku=0$ 有解 $u\neq 0$。
\end{enumerate}
此二择一是 Fredholm 二择一。此外，若 ($\beta$) 成立，齐次问题的解空间是有限维的，非齐次方程
\begin{align}
u-Ku=f
\label{align:appD-Fredholm非齐次}
\end{align}
有解当且仅当 $f\in N(I-K^*)^\perp$。
\end{remark}

现研究紧线性算子的谱。

\begin{theorem}[紧算子的谱。]\label{theorem:appD-紧算子谱}假设 $\dim H=\infty$ 且 $K:H\to H$ 紧。则
\begin{enumerate}[(i)]
\item $0\in\sigma(K)$，
\item $\sigma(K)\setminus\{0\}=\sigma_p(K)\setminus\{0\}$，
\item $\sigma(K)\setminus\{0\}$ 有限，否则 $\sigma(K)\setminus\{0\}$ 是趋于 $0$ 的序列。
\end{enumerate}
\end{theorem}

\begin{proof}
1. 假设 $0\notin\sigma(K)$。则 $K:H\to H$ 双射，故 $I=K\circ K^{-1}$ 作为紧算子和有界线性算子的复合是紧的。这不可能，因 $\dim H=\infty$。

2. 假设 $\eta\in\sigma(K)$，$\eta\neq 0$。则若 $N(K-\eta I)=\{0\}$，Fredholm 二择一将蕴含 $R(K-\eta I)=H$。但此时 $\eta\in\rho(K)$，矛盾。

3. 现假设 $\{\eta_k\}_{k=1}^\infty$ 是 $\sigma(K)\setminus\{0\}$ 的不同元素序列且 $\eta_k\to\eta$。将证明 $\eta=0$。

事实上，因 $\eta_k\in\sigma_p(K)$，存在 $w_k\neq 0$ 使得 $Kw_k=\eta_kw_k$。设 $H_k$ 表示 $H$ 中由 $\{w_1,\cdots,w_k\}$ 张成的子空间。则 $H_k\subsetneq H_{k+1}$ 对每个 $k=1,2,\cdots$，因 $\{w_k\}_{k=1}^\infty$ 线性无关。还注意 $(K-\eta_kI)H_k\subseteq H_{k-1}$（$k=2,\cdots$）。现对 $k=1,\cdots$ 选取元素 $u_k\in H_k$，$u_k\in H_{k-1}^\perp$，$\|u_k\|=1$。若 $k>l$，$H_{l-1}\subsetneq H_l\subseteq H_{k-1}\subsetneq H_k$。因此
\begin{align*}
\left\|\frac{Ku_k}{\eta_k}-\frac{Ku_l}{\eta_l}\right\|=\left\|\frac{Ku_k-\eta_ku_k}{\eta_k}-\frac{Ku_l-\eta_lu_l}{\eta_l}+u_k-u_l\right\|\geqslant 1,
\end{align*}
因 $Ku_k-\eta_ku_k$，$Ku_l-\eta_lu_l$，$u_l\in H_{k-1}$。若 $\eta_k\to\eta\neq 0$，我们得到与 $K$ 紧性的矛盾。
\end{proof}

\section{对称算子}\label{section:appD-对称算子}

现设 $S:H\to H$ 线性、有界、对称，写
\begin{align*}
m\coloneqq\inf_{\substack{u\in H\\\|u\|=1}}(Su,u),\quad M\coloneqq\sup_{\substack{u\in H\\\|u\|=1}}(Su,u).
\end{align*}

\begin{lemma}[谱的界。]\label{lemma:appD-谱界}我们有
\begin{enumerate}[(i)]
\item $\sigma(S)\subset[m,M]$，
\item $m,M\in\sigma(S)$。
\end{enumerate}
\end{lemma}

\begin{proof}
1. 设 $\eta>M$。则
\begin{align*}
(\eta u-Su,u)\geqslant(\eta-M)\|u\|^2 \quad (u\in H).
\end{align*}
因此 Lax--Milgram 定理（\S6.2.1）断言 $\eta I-S$ 是一一对应且到上的，故 $\eta\in\rho(S)$。类似地若 $\eta<m$ 则 $\eta\in\rho(S)$。这证明了 (i)。

2. 将证明 $M\in\sigma(S)$。因配对 $[u,v]\coloneqq(Mu-Su,v)$ 是对称的，$[u,u]\geqslant 0$ 对所有 $u\in H$，Cauchy--Schwarz 不等式蕴含
\begin{align*}
|(Mu-Su,v)|\leqslant(Mu-Su,u)^{1/2}(Mv-Sv,v)^{1/2}
\end{align*}
对所有 $u,v\in H$。特别地
\begin{align}
\|Mu-Su\|\leqslant C(Mu-Su,u)^{1/2} \quad (u\in H)
\label{align:appD-对称估计}
\end{align}
对某常数 $C$。

现设 $\{u_k\}_{k=1}^\infty\subset H$ 满足 $\|u_k\|=1$（$k=1,\cdots$）且 $(Su_k,u_k)\to M$。则\eqref{align:appD-对称估计} 蕴含 $\|Mu_k-Su_k\|\to 0$。现若 $M\in\rho(S)$，则
\begin{align*}
u_k=(MI-S)^{-1}(Mu_k-Su_k)\to 0,
\end{align*}
矛盾。因此 $M\in\sigma(S)$，类似地 $m\in\sigma(S)$。
\end{proof}

\begin{theorem}[紧对称算子的特征向量。]\label{theorem:appD-紧对称特征向量}设 $H$ 是可分 Hilbert 空间，假设 $S:H\to H$ 是紧对称算子。则存在 $H$ 的由 $S$ 的特征向量组成的可数正交归一基。
\end{theorem}

\begin{proof}
1. 设 $\{\eta_k\}$ 构成 $S$ 的除 $0$ 外的不同特征值序列。设 $\eta_0=0$。写 $H_0=N(S)$，$H_k=N(S-\eta_kI)$（$k=1,\cdots$）。则 $0\leqslant\dim H_0\leqslant\infty$，$0<\dim H_k<\infty$，由 Fredholm 二择一。

2. 设 $u\in H_k$，$v\in H_l$（$k\neq l$）。则 $Su=\eta_ku$，$Sv=\eta_lv$，故
\begin{align*}
\eta_k(u,v)=(Su,v)=(u,Sv)=\eta_l(u,v).
\end{align*}
因 $\eta_k\neq\eta_l$，推出 $(u,v)=0$。因此子空间 $H_k$ 和 $H_l$ 正交。

3. 现设 $\widetilde{H}$ 是 $H$ 中包含 $H_0,H_1,\cdots$ 的最小子空间。于是
\begin{align*}
\widetilde{H}=\left\{\sum_{k=0}^m a_ku_k\mid m\in\{0,\cdots\},\;u_k\in H_k,\;a_k\in\mathbb{R}\right\}.
\end{align*}
接下来证明 $\widetilde{H}$ 在 $H$ 中稠密。显然 $S(\widetilde{H})\subseteq\widetilde{H}$。此外 $S(\widetilde{H}^\perp)\subseteq\widetilde{H}^\perp$：事实上若 $u\in\widetilde{H}^\perp$ 且 $v\in\widetilde{H}$，则 $(Su,v)=(u,Sv)=0$。

现算子 $\widetilde{S}\equiv S|_{\widetilde{H}^\perp}$ 是紧对称的。此外 $\sigma(\widetilde{S})=\{0\}$，因 $\widetilde{S}$ 的任何非零特征值也是 $S$ 的特征值。由引理，$(\widetilde{S}u,u)=0$ 对所有 $u\in\widetilde{H}^\perp$。但若 $u,v\in\widetilde{H}^\perp$，
\begin{align*}
2(\widetilde{S}u,v)=(\widetilde{S}(u+v),u+v)-(\widetilde{S}u,u)-(\widetilde{S}v,v)=0.
\end{align*}
故 $\widetilde{S}=0$。因此 $\widetilde{H}^\perp\subset N(S)\subset\widetilde{H}$，故 $\widetilde{H}^\perp=\{0\}$。因此 $\widetilde{H}$ 在 $H$ 中稠密。

4. 对每个子空间 $H_k$（$k=0,\cdots$）选取正交归一基，注意因 $H$ 可分，$H_0$ 有可数正交归一基。由此得到特征向量的正交归一基。
\end{proof}

这些证明大多来自 Brezis~\cite{Br1}。另见 Gilbarg--Trudinger~\cite{G-T} 第 5 章、Lax~\cite{Lx2}、Reed--Simon~\cite{R-S1} 和 Yosida~\cite{Y}。

\end{document}